\documentclass[a4paper,11pt,oneside]{article}


%%%% Packages: %%%%%%%%%%%%%%%%%%%%%%%%%%%%%%%%%%

\usepackage[utf8]{inputenc}
\usepackage[T1]{fontenc}

% \usepackage[english]{babel} 
\usepackage[english]{babel} 

\usepackage{cite}

\usepackage{geometry}
\geometry{textwidth=16cm,textheight=26cm}

%%%% Preamble: %%%%%%%%%%%%%%%%%%%%%%%%%%%%%%%%%%

\title{D2. Bibliography : Design of Complex CI Pipelines}

\author{
    Brice Delcroix
    \and Timoth\'{e} Piscaglia
    \and Pierre-Louis Lef\`{e}vre}
    
\date{
    Vint Cerf, INFOB302 - IDS (2025 - 2026)
}
    
%%%% Document: %%%%%%%%%%%%%%%%%%%%%%%%%%%%%%%%%%

\begin{document}

\maketitle

% ---------------------------------------
\section{Search criteria}
% ---------------------------------------

This document includes a list of published \textbf{peer-reviewed} research papers collected following a search in bibliographic engines or via \textit{snowballing}. The search queries used are:
%
\begin{itemize}
    \item \textbf{Scopus:} ``continuous integration'' AND ``practice'' AND ``modeling'' AND ``industry''
    \item \textbf{Scopus:} ``continuous integration'' AND ``configuration'' AND ``anti-pattern'' AND ``Travis CI''
    \item \textbf{Scopus:} ``continuous delivery'' AND ``configuration smells'' AND ``pipeline'' AND ``linter''
    \item \textbf{Scopus:} ``CI/CD pipeline'' AND ``evolution'' AND ``restructuring''
    \item \textbf{Scopus:} ``continuous integration'' AND ``bad practices'' AND ``empirical''
    \item \textbf{Google Scholar:} ``GitHub Actions'' AND ``workflow'' AND ``continuous integration'' AND ``empirical study''
    \item \textbf{Scopus:} ``architecting'' AND ``continuous delivery'' AND ``deployment'' AND ``empirical''
    \item \textbf{Google Scholar:} ``code review'' AND ``build system'' AND ``specifications'' AND ``continuous integration''
    \item \textbf{Scopus:} ``continuous integration'' AND ``anti-patterns'' AND ``decay'' AND ``automated''
    \item \textbf{Google Scholar:} ``continuous integration'' AND ``operational data'' AND ``CircleCI'' AND ``case study''
\end{itemize}

% ---------------------------------------
\section{List of papers}
% ---------------------------------------

% ---- Paper 1 ----
\subsection{Modeling Continuous Integration Practice Differences in Industry Software Development \cite{Stahl2014}}
%
\begin{itemize}
    \item \textbf{Found:} using Scopus with query ``continuous integration'' AND ``practice'' AND ``modeling'' AND ``industry''.
    \item \textbf{Justification:} this paper proposes a descriptive model to characterize variation points in how CI practices are implemented across industry projects, including scope, trigger mechanisms, and integration frequency. It provides a foundational vocabulary for reasoning about CI pipeline architecture and design diversity, which is directly relevant to understanding how complex CI pipelines are designed.
    \item \textbf{Cross-validation:} found by \ldots{} confirmed/rejected by \ldots{}
\end{itemize}

% ---- Paper 2 ----
\subsection{Use and Misuse of Continuous Integration Features: An Empirical Study of Projects That (Mis)Use Travis CI \cite{Gallaba2020}}
%
\begin{itemize}
    \item \textbf{Found:} using Scopus with query ``continuous integration'' AND ``configuration'' AND ``anti-pattern'' AND ``Travis CI''.
    \item \textbf{Justification:} this landmark study analyzes CI configuration feature usage and misuse across 9,312 open-source Travis CI projects. It identifies four anti-patterns in CI specification files and reveals that 48\% of configuration code is devoted to job environment setup. The paper proposes Gretel, a tool to detect CI configuration anti-patterns at scale, directly addressing configuration management challenges in complex pipelines.
    \item \textbf{Cross-validation:} found by \ldots{} confirmed/rejected by \ldots{}
\end{itemize}

% ---- Paper 3 ----
\subsection{Configuration Smells in Continuous Delivery Pipelines: A Linter and a Six-Month Study on GitLab \cite{Vassallo2020}}
%
\begin{itemize}
    \item \textbf{Found:} using Scopus with query ``continuous delivery'' AND ``configuration smells'' AND ``pipeline'' AND ``linter''.
    \item \textbf{Justification:} this paper proposes CD-Linter, a semantic linter that detects four types of configuration smells in pipeline YAML files. Evaluated on 5,312 open-source GitLab projects over six months, it achieves 87\% precision and 94\% recall. The study is central to the question of pipeline configuration maintainability in complex CI/CD environments.
    \item \textbf{Cross-validation:} found by \ldots{} confirmed/rejected by \ldots{}
\end{itemize}

% ---- Paper 4 ----
\subsection{CI/CD Pipelines Evolution and Restructuring: A Qualitative and Quantitative Study \cite{Zampetti2021}}
%
\begin{itemize}
    \item \textbf{Found:} using Scopus with query ``CI/CD pipeline'' AND ``evolution'' AND ``restructuring''.
    \item \textbf{Justification:} this paper combines qualitative analysis of 615 pipeline configuration change commits with quantitative analysis of 253,492 commits from 4,644 Travis CI projects. It produces a taxonomy of 34 CI/CD pipeline restructuring actions and 16 evolution metrics, directly characterizing how pipeline components evolve and accumulate technical debt over time.
    \item \textbf{Cross-validation:} found by \ldots{} confirmed/rejected by \ldots{}
\end{itemize}

% ---- Paper 5 ----
\subsection{An Empirical Characterization of Bad Practices in Continuous Integration \cite{Zampetti2020}}
%
\begin{itemize}
    \item \textbf{Found:} using Scopus with query ``continuous integration'' AND ``bad practices'' AND ``empirical''.
    \item \textbf{Justification:} this study empirically investigates bad practices in CI through developer interviews, Stack Overflow analysis, and online surveys. It compiles a comprehensive catalog of 79 CI bad smells across 7 categories, providing the most extensive taxonomy of maintenance-related anti-patterns in CI configurations. This is directly relevant to understanding what makes CI pipelines complex and error-prone.
    \item \textbf{Cross-validation:} found by \ldots{} confirmed/rejected by \ldots{}
\end{itemize}

% ---- Paper 6 ----
\subsection{How Do Software Developers Use GitHub Actions to Automate Their Workflows? \cite{Kinsman2021}}
%
\begin{itemize}
    \item \textbf{Found:} using Google Scholar with query ``GitHub Actions'' AND ``workflow'' AND ``continuous integration'' AND ``empirical study''.
    \item \textbf{Justification:} this is the first study investigating how developers design and use GitHub Actions YAML-based workflow files to automate CI/CD pipelines. It categorizes actions by purpose (CI, deployment, code quality) and examines how pipeline-as-code configurations are structured in practice, offering direct insight into modern CI workflow design patterns.
    \item \textbf{Cross-validation:} found by \ldots{} confirmed/rejected by \ldots{}
\end{itemize}

% ---- Paper 7 ----
\subsection{An Empirical Study of Architecting for Continuous Delivery and Deployment \cite{Shahin2019}}
%
\begin{itemize}
    \item \textbf{Found:} using Scopus with query ``architecting'' AND ``continuous delivery'' AND ``deployment'' AND ``empirical''.
    \item \textbf{Justification:} this paper presents a conceptual framework for (re-)architecting software systems to support CD pipelines, based on 21 practitioner interviews and a 91-respondent survey. It directly addresses how software architecture decisions shape pipeline design, including deployment unit granularity and service decomposition, which are key concerns when designing complex CI pipelines.
    \item \textbf{Cross-validation:} found by \ldots{} confirmed/rejected by \ldots{}
\end{itemize}

% ---- Paper 8 ----
\subsection{Code Review of Build System Specifications: Prevalence, Purposes, Patterns, and Perceptions \cite{Nejati2023}}
%
\begin{itemize}
    \item \textbf{Found:} using Google Scholar with query ``code review'' AND ``build system'' AND ``specifications'' AND ``continuous integration''.
    \item \textbf{Justification:} this is the first empirical study of code review practices for build system specifications (pipeline-as-code). Analyzing 502,931 change sets from Qt and Eclipse, it finds that build specs are reviewed two times less frequently than production code, yet review comments more often identify defects. This is directly relevant to assessing the design quality of complex pipeline specifications.
    \item \textbf{Cross-validation:} found by \ldots{} confirmed/rejected by \ldots{}
\end{itemize}

% ---- Paper 9 ----
\subsection{Automated Reporting of Anti-Patterns and Decay in Continuous Integration \cite{Vassallo2019}}
%
\begin{itemize}
    \item \textbf{Found:} using Scopus with query ``continuous integration'' AND ``anti-patterns'' AND ``decay'' AND ``automated''.
    \item \textbf{Justification:} this paper proposes CI-Odor, an automated tool detecting four key CI anti-patterns (Slow Build, Broken Release, Late Merging, Missing Build Failure Repair) by analyzing build logs and repository data. It demonstrates that automated detection can identify CI pipeline decay early, directly targeting the maintenance dimension of complex CI pipelines.
    \item \textbf{Cross-validation:} found by \ldots{} confirmed/rejected by \ldots{}
\end{itemize}

% ---- Paper 10 ----
\subsection{Lessons from Eight Years of Operational Data from a Continuous Integration Service: An Exploratory Case Study of CircleCI \cite{Gallaba2022}}
%
\begin{itemize}
    \item \textbf{Found:} using Google Scholar with query ``continuous integration'' AND ``operational data'' AND ``CircleCI'' AND ``case study''.
    \item \textbf{Justification:} this large-scale empirical study uses eight years of operational data from CircleCI, a major CI service provider. It examines CI service scalability challenges, configuration patterns, and usage trends across thousands of projects, providing a unique service-provider perspective on how CI pipelines scale and operate in production environments.
    \item \textbf{Cross-validation:} found by \ldots{} confirmed/rejected by \ldots{}
\end{itemize}

% Do not forget to add the publications to the bib file to get the full reference.
\bibliographystyle{plain}
\bibliography{biblio}

\end{document}
